\chapter{\IfLanguageName{dutch}{Conclusie}{Conclusion}}%
\label{ch:conclusie}

Dit onderzoek vertrok vanuit een concrete en maatschappelijk relevante
probleemstelling: bezoekers van begraafplaatsen ervaren moeilijkheden bij het
terugvinden van een specifiek graf, terwijl bestaande digitale oplossingen nog
te veel vragen dat de gebruiker zijn aandacht afwisselt tussen scherm en omgeving.
De centrale onderzoeksvraag luidde: \textit{Op welke manier kan digitalisatie,
    in combinatie met Augmented Reality, worden ingezet om de efficiëntie van het
    navigeren op een kerkhof te verbeteren?}

\section{Antwoord op de centrale onderzoeksvraag}

Op basis van dit onderzoek kan worden besloten dat augmented reality een
technisch haalbare en functionele oplossing biedt voor kerkhofnavigatie. De
ontwikkelde proof of concept toont aan dat een 3D AR-navigatiepijl in combinatie
met een world-locked grafmarker de gebruiker correct en intuïtief naar een
geselecteerd graf kan leiden, zonder dat de gebruiker zijn blik op de omgeving
moet veranderen. De technische evaluatie bevestigde dat alle kerncomponenten correct
functioneren: de pijl wees gedurende de volledige test in de juiste richting,
de grafmarker werd zichtbaar op circa 20 meter, de ARKit world tracking was
stabiel vanaf de eerste seconde, en de zoekfunctie, minikaart en afstandslabel
functioneerden allen correct.

Tegelijkertijd maakt dit onderzoek duidelijk dat AR geen op zichzelf staande
oplossing is. De combinatie met een digitale kaartweergave is essentieel: de
minikaart biedt de ruimtelijke context die pure AR mist, en fungeert als
betrouwbare fallback wanneer AR-tracking tijdelijk instabiel is. Dit hybride
model, waarbij AR en kaart elkaar aanvullen in plaats van vervangen, is de
sterkste bevinding van dit onderzoek.

\section{Technische bevindingen}

\begin{itemize}
    \item \textbf{ARKit .gravityAndHeading} maakt directe GPS-naar-AR conversie
    mogelijk via een vlakke-Aarde-benadering met een fout van minder dan 0,01\%
    binnen 500 meter, wat zeker volstaat voor kerkhofschaal.
    
    \item \textbf{GPS-nauwkeurigheid} blijft de voornaamste beperking.
    Op een begraafplaats met grafstenen die slechts 60 tot 80 centimeter uit
    elkaar staan, is een GPS-nauwkeurigheid van 3 tot 10 meter onvoldoende voor
    exacte grafsteen-identificatie. Duidelijke visuele markers en officieel
    ingemeten coördinaten zijn noodzakelijk om dit te compenseren.
    
    \item \textbf{Padsgewijze navigatie} ontbreekt. De huidige pijlnavigatie
    wijst steeds in een rechte lijn naar het graf. Voor een optimale
    gebruikerservaring is een routelijn langs de bestaande paden gewenst,
    waarvoor GIS-paddata vereist is. OpenStreetMap biedt hiervoor een
    schaalbare oplossing \autocite{Novack2018}.
    
    \item \textbf{Platformonafhankelijkheid} is aangetoond: alle technische
    componenten hebben een directe equivalent op Android via ARCore en Google's
    Geospatial API. Een cross-platform implementatie via Unity AR Foundation
    is technisch haalbaar \autocite{GoogleARCore2024}.
\end{itemize}

\section{Juridische bevindingen}

\begin{itemize}
    \item Grafdata van overledenen valt buiten de bescherming van de GDPR,
    maar vereist zorgvuldig gebruik vanwege de mogelijke indirecte link met
    levende familieleden.
    \item Het rijksregisternummer mag onder geen enkele reden worden verwerkt in een
    publieke navigatieapplicatie op basis van de Wet van 8 augustus 1983
    \autocite{WetRijksregister1983}.
    \item Voor een productieomgeving zijn een verwerkersovereenkomst met de
    gemeente, een privacyverklaring en mogelijk een DPIA vereiste stappen.
\end{itemize}

\section{Beperkingen van dit onderzoek}

//todo
%todo aanvullen

\section{Maatschappelijke relevantie}

Dit onderzoek toont aan dat technologie, ook in een emotioneel gevoelige context
zoals een begraafplaats, een zinvolle bijdrage kan leveren aan de gebruikerservaring.
Augmented reality is hierbij geen doel op zich, maar een middel om een concreet
probleem op te lossen: het gemakkelijker en minder stressvol maken van het
terugvinden van een graf. De keuze voor een discreet, snel en intuïtief ontwerp
weerspiegelt het idee dat technologie in deze context de achtergrond moet
blijven, en de mens op de voorgrond.

Indien deze proof of concept verder wordt ontwikkeld en geïmplementeerd door de
gemeente Sint-Gillis-Waas of een vergelijkbare gemeente, betekent het een
concrete verbetering voor elke bezoeker die op zoek is naar een dierbare. Dat
is uiteindelijk de kern van dit onderzoek.